\PassOptionsToPackage{unicode=true}{hyperref} % options for packages loaded elsewhere
\PassOptionsToPackage{hyphens}{url}
%
\documentclass[14pt,]{extarticle}
\usepackage{lmodern}
\usepackage{amssymb,amsmath}
\usepackage{ifxetex,ifluatex}
\usepackage{fixltx2e} % provides \textsubscript
\ifnum 0\ifxetex 1\fi\ifluatex 1\fi=0 % if pdftex
  \usepackage[T1]{fontenc}
  \usepackage[utf8]{inputenc}
  \usepackage{textcomp} % provides euro and other symbols
\else % if luatex or xelatex
  \usepackage{unicode-math}
  \defaultfontfeatures{Ligatures=TeX,Scale=MatchLowercase}
    \setmainfont[]{Times New Roman}
  \ifxetex
    \usepackage{xeCJK}
    \setCJKmainfont[]{Microsoft JhengHei}
  \fi
  \ifluatex
    \usepackage[]{luatexja-fontspec}
    \setmainjfont[]{Microsoft JhengHei}
  \fi
\fi
% use upquote if available, for straight quotes in verbatim environments
\IfFileExists{upquote.sty}{\usepackage{upquote}}{}
% use microtype if available
\IfFileExists{microtype.sty}{%
\usepackage[]{microtype}
\UseMicrotypeSet[protrusion]{basicmath} % disable protrusion for tt fonts
}{}
\IfFileExists{parskip.sty}{%
\usepackage{parskip}
}{% else
\setlength{\parindent}{0pt}
\setlength{\parskip}{6pt plus 2pt minus 1pt}
}
\usepackage{hyperref}
\hypersetup{
            pdfborder={0 0 0},
            breaklinks=true}
\urlstyle{same}  % don't use monospace font for urls
\usepackage[margin=1in]{geometry}
\usepackage{longtable,booktabs}
% Fix footnotes in tables (requires footnote package)
\IfFileExists{footnote.sty}{\usepackage{footnote}\makesavenoteenv{longtable}}{}
\setlength{\emergencystretch}{3em}  % prevent overfull lines
\providecommand{\tightlist}{%
  \setlength{\itemsep}{0pt}\setlength{\parskip}{0pt}}
\setcounter{secnumdepth}{0}
% Redefines (sub)paragraphs to behave more like sections
\ifx\paragraph\undefined\else
\let\oldparagraph\paragraph
\renewcommand{\paragraph}[1]{\oldparagraph{#1}\mbox{}}
\fi
\ifx\subparagraph\undefined\else
\let\oldsubparagraph\subparagraph
\renewcommand{\subparagraph}[1]{\oldsubparagraph{#1}\mbox{}}
\fi

% set default figure placement to htbp
\makeatletter
\def\fps@figure{htbp}
\makeatother


\date{}

\begin{document}

\hypertarget{week-06-ux4f5cux696dux56deux994bposner-spatial-cueing-task-builder-pavlovia}{%
\section{Week 06 作業回饋（Posner Spatial Cueing Task --- Builder +
Pavlovia）}\label{week-06-ux4f5cux696dux56deux994bposner-spatial-cueing-task-builder-pavlovia}}

\hypertarget{ux4f55ux5b98ux81fb-92100}{%
\subsection{114825002 何官臻 ---
92／100}\label{ux4f55ux5b98ux81fb-92100}}

\hypertarget{ux7e3dux8a55}{%
\subsubsection{總評}\label{ux7e3dux8a55}}

\textbf{Flow 結構完整度最高}，嚴格依照作業指定的 8 個
routine（Instructions → Fixation → Cue → SOA\_Blank → Target → ITI →
Rest → Goodbye）實作，並使用 nested loop（\texttt{block\_loop} ×
\texttt{trial\_loop}）實現「2 blocks × 20 trials + Rest」的區塊設計
✓。Cue 與 Target routine 皆有 CodeComponent 控制位置，且已完成 Pavlovia
線上部署（\texttt{Pavlovia\_data\_*.csv} 即為線上執行產出，證明 JS
版本可實際運作）。扣分主要來自條件檔 \texttt{validity}
欄位的標籤錯誤，以及 \texttt{readme.md} 空白缺少文件說明。

\hypertarget{ux5404ux9805ux8a55ux8a9e}{%
\subsubsection{各項評語}\label{ux5404ux9805ux8a55ux8a9e}}

\begin{longtable}[]{@{}lll@{}}
\toprule
\begin{minipage}[b]{0.09\columnwidth}\raggedright
指標\strut
\end{minipage} & \begin{minipage}[b]{0.09\columnwidth}\raggedright
得分\strut
\end{minipage} & \begin{minipage}[b]{0.09\columnwidth}\raggedright
評語\strut
\end{minipage}\tabularnewline
\midrule
\endhead
\begin{minipage}[t]{0.09\columnwidth}\raggedright
Experiment Structure\strut
\end{minipage} & \begin{minipage}[t]{0.09\columnwidth}\raggedright
20/20\strut
\end{minipage} & \begin{minipage}[t]{0.09\columnwidth}\raggedright
8 個 routine
全部到位：Instructons、Fixation、Cue、SOA\_Blank、Target、ITI、Rest、Goodbye，Flow
排序正確，為本週作業中最完整的結構實作。\strut
\end{minipage}\tabularnewline
\begin{minipage}[t]{0.09\columnwidth}\raggedright
Conditions File\strut
\end{minipage} & \begin{minipage}[t]{0.09\columnwidth}\raggedright
12/15\strut
\end{minipage} & \begin{minipage}[t]{0.09\columnwidth}\raggedright
已包含全部 5 個必要欄位
\texttt{cue\_side,\ target\_side,\ validity,\ soa,\ correct\_key} ✓，8
valid + 2 invalid 比例正確 ✓。但 \texttt{validity} 欄位中「valid」的 8
列皆誤填為 \texttt{"validity"}（應為
\texttt{"valid"}），會造成後續資料分析依 validity
分組時出錯。\textbf{(-3)}\strut
\end{minipage}\tabularnewline
\begin{minipage}[t]{0.09\columnwidth}\raggedright
Loop \& Randomization\strut
\end{minipage} & \begin{minipage}[t]{0.09\columnwidth}\raggedright
10/10\strut
\end{minipage} & \begin{minipage}[t]{0.09\columnwidth}\raggedright
\texttt{block\_loop} nReps=2 外包 \texttt{trial\_loop} nReps=2 × 10 列 =
40 試驗，兩層皆設為 random ✓，Rest 在 trial\_loop 外、block\_loop
內，位置正確。\strut
\end{minipage}\tabularnewline
\begin{minipage}[t]{0.09\columnwidth}\raggedright
Cue \& Target Positioning\strut
\end{minipage} & \begin{minipage}[t]{0.09\columnwidth}\raggedright
15/15\strut
\end{minipage} & \begin{minipage}[t]{0.09\columnwidth}\raggedright
Cue 與 Target routine 各自包含 CodeComponent 依
\texttt{\$cue\_side}／\texttt{\$target\_side} 設定位置；Polygon 與 Text
的 Position 欄位使用 set every repeat，設計正確。\strut
\end{minipage}\tabularnewline
\begin{minipage}[t]{0.09\columnwidth}\raggedright
Keyboard Response\strut
\end{minipage} & \begin{minipage}[t]{0.09\columnwidth}\raggedright
10/10\strut
\end{minipage} & \begin{minipage}[t]{0.09\columnwidth}\raggedright
Target routine 中 \texttt{key\_resp} 限定左右方向鍵、store correct、最大
1.5 秒 ✓，CSV 中 \texttt{key\_resp.corr}／\texttt{key\_resp.rt}
正確記錄。\strut
\end{minipage}\tabularnewline
\begin{minipage}[t]{0.09\columnwidth}\raggedright
Pavlovia Deployment\strut
\end{minipage} & \begin{minipage}[t]{0.09\columnwidth}\raggedright
15/15\strut
\end{minipage} & \begin{minipage}[t]{0.09\columnwidth}\raggedright
已提供連結
\texttt{https://pavlovia.org/ireneho3507/posner\_task}，並附上實際線上執行的
\texttt{Pavlovia\_data\_week06\_hw\_2026-04-07\_23h58.25.416.csv}，證明
JS 版本可正常運作。\strut
\end{minipage}\tabularnewline
\begin{minipage}[t]{0.09\columnwidth}\raggedright
Pilot Data\strut
\end{minipage} & \begin{minipage}[t]{0.09\columnwidth}\raggedright
10/10\strut
\end{minipage} & \begin{minipage}[t]{0.09\columnwidth}\raggedright
Pavlovia CSV 欄位齊全（含
\texttt{key\_resp.corr}、\texttt{key\_resp.rt}、\texttt{cue\_side}、\texttt{target\_side}、\texttt{trial\_loop.thisN}
等），另附本機執行的 \texttt{local\_run\_data}
資料夾作為備份，資料品質佳。\strut
\end{minipage}\tabularnewline
\begin{minipage}[t]{0.09\columnwidth}\raggedright
Documentation\strut
\end{minipage} & \begin{minipage}[t]{0.09\columnwidth}\raggedright
0/5\strut
\end{minipage} & \begin{minipage}[t]{0.09\columnwidth}\raggedright
⚠️ \texttt{readme.md} 為\textbf{空檔案}（0 行），未提供設計說明或 Code
Component 用途描述，不符合作業「Brief description explaining your design
choices」要求。\textbf{(-5)}\strut
\end{minipage}\tabularnewline
\bottomrule
\end{longtable}

\hypertarget{ux6539ux9032ux5efaux8b70}{%
\subsubsection{改進建議}\label{ux6539ux9032ux5efaux8b70}}

\begin{enumerate}
\def\labelenumi{\arabic{enumi}.}
\tightlist
\item
  \textbf{修正 \texttt{validity} 欄位標籤（+3 分）：} 將 8 列有效試驗的
  \texttt{validity} 從 \texttt{"validity"} 改為
  \texttt{"valid"}，否則統計時 \texttt{df{[}df.validity=="valid"{]}}
  將回傳 0 筆。
\item
  \textbf{補寫 README（+5 分）：} 說明：(a) 為何選擇 nested block\_loop
  / trial\_loop 結構；(b) Cue 與 Target 兩處 CodeComponent 的 Python↔JS
  對應；(c) CSV 欄位含意。完成後可拿回全部 5 分。
\end{enumerate}

\end{document}
